\chapter{Evaluation protocol}
\label{cha:protocol}

The dataset contract in Section~\ref{sec:contract} defines the forecast sample. This chapter specifies how those samples are compared. It defines which plants are held out, what each model may observe at the forecast origin, how errors are aggregated, and how visual reliance is tested. Within each dataset track, direct comparisons use the same forecast origins, target masks, and horizon.

\section{Plant-disjoint partitions}

Plant identity is the unit of partitioning within each dataset track. The fixed United Kingdom assignment contains 69 training plants, 15 validation plants, and 14 test plants. Table~\ref{tab:uk-partition} gives the quality-valid power observations in each group before forecast-window and daylight filtering. Any evaluation on another track requires its own plant-disjoint partition under the same separation rule. The results chapter reports the UK test partition.

\begin{table}[htbp]
\centering
\begin{tabular}{@{}lccl@{}}
\hline
Split & Plants & Valid steps & Use \\
\hline
Training & 69 & 846,633 & Fitting and adaptation \\
Validation & 15 & 182,809 & Early stopping and tuning \\
Test & 14 & 171,543 & Final evaluation \\
\hline
\end{tabular}
\caption{United Kingdom plant partition after quality-controlled exclusions. Each installation belongs to exactly one partition.}
\label{tab:uk-partition}
\end{table}

The evaluator advances by one forecast horizon between origins and applies the target-validity and daylight masks to the windows defined in Section~\ref{sec:contract}. For the United Kingdom partition, this leaves 19,964 test windows and 165,295 scored target steps. The 171,543 test steps in Table~\ref{tab:uk-partition} count valid power observations before those filters.

Forecast origins denote the timestamps separating historical observations from future horizon steps. Eligible evaluation points are governed by two filters: a target-validity mask excluding sensor faults and missing intervals, and a daylight mask restricting evaluation to solar elevation angles above zero. Evaluating all models across identical masked steps guarantees that performance differences reflect true forecasting quality under uniform test periods.

Training plants supply fitted parameters, including adapter weights, and any retrieval datastore. Validation plants determine checkpoints and hyperparameters, including retrieval weights when tuned. Test plants are used only for final reporting. This boundary also applies to learned normalization and feature selection. Each implementation declares its fitted statistics. The loader checks that plant assignments do not overlap.

\section{Scenarios and input parity}

The primary benchmark is the United Kingdom plant-disjoint comparison. The cross-region experiment applies models trained on UK plants to all eight eligible US plants. Checkpoints from seeds 42, 43, and 44 remain fixed during US inference, and US power is resampled to 30-minute cadence. Thus, the two reported settings measure transfer to unseen plants within the UK and transfer from the UK to the US. Across these comparisons, each model receives the modalities its design supports. Persistence forecasts use no satellite frames. Numerical forecasters use power history and, where supported, covariates. Multimodal systems can also use recent imagery. Table~\ref{tab:input-parity} summarizes the inputs by model family. Each result identifies the fields used by the individual method.

\begin{table}[htbp]
\centering
\begin{tabularx}{\textwidth}{@{}lcccY@{}}
\hline
Family & Power & Covariates & Images & Role \\
\hline
References & Yes & Varies & No & Daily/short-horizon reference \\
Supervised models & Yes & Varies & No & Numerical forecasting \\
Tabular FMs & Yes & Yes & No & Tabular pretraining \\
Time-series FMs & Yes & Varies & No & Time-series pretraining \\
Retrieval/adapters & Yes & Varies & No & Retrieval or adaptation \\
Multimodal models & Yes & Varies & Yes & Visual context \\
\hline
\end{tabularx}
\caption{Modalities available at the shared forecast origin. FMs denotes foundation models; Varies means method-specific covariates, where supported.}
\label{tab:input-parity}
\end{table}

\section{Baseline implementations and comparison boundaries}
\label{sec:baseline-implementations}

The models reviewed in Chapter~\ref{cha:related} are compared under the track-specific plant split and input rules above. This subsection specifies each model’s inputs, whether it learns from the training plants, and whether its reported score uses the same evaluation set as other methods.

The non-learned references use only permitted observations. Persistence repeats the latest valid power value across the horizon. Seasonal naive reuses the value from the same clock time on the previous day. Hourly climatology estimates a plant's typical output at each hour from its available history. Smart Persistence carries the latest observed clearness relationship forward as the deterministic solar envelope changes~\cite{haurwitz}. It supplies the denominator of the skill score.

Tabular models (LightGBM~\cite{lightgbm} and TabPFN~\cite{tabpfn3}) operate on fixed feature tables of lagged power, calendar variables, and weather observations without satellite frames or future covariates. Supervised sequence baselines (PatchTST~\cite{patchtst}, TFT~\cite{tft}, DLinear~\cite{dlinear}, and an MLP baseline) train on the same plant partition using historical covariates. iTransformer~\cite{itransformer} models cross-variable interactions directly, providing a non-visual multi-variable reference.

The time-series foundation-model entries start from released checkpoints. Chronos-2~\cite{chronos2} is evaluated both zero-shot and after adaptation on training plants as the numerical control of the fusion ladder. TimesFM~\cite{timesfm} and TiRex~\cite{tirex} are used zero-shot. Tiny Time Mixer~\cite{ttm} appears in zero-shot and fine-tuned forms. CoRA~\cite{cora} adds a covariate adapter to a mostly frozen backbone. TS-RAG~\cite{tsrag} and Cross-RAG~\cite{crossrag} condition the forecaster on windows retrieved from training plants. These entries receive no imagery. Their architectures, pretraining corpora, and adaptation strategies differ. Their rankings therefore compare benchmark performance across methods and do not constitute an architectural ablation.

Multimodal baselines encompass both satellite-conditioned models (MMTSFM, Aurora, Solar-VLM, SUNSET) and visual time-series rendering methods (Time-VLM, UniCast). All models are scored across the held-out test installations on valid daylight steps.

\section{Scoring and reporting}

Within each track, valid daylight targets are scored separately for every held-out plant before the plant scores are averaged. This gives installations equal weight even when their usable records differ in length. The two co-primary measures answer different questions. Skill score measures the reduction in NRMSE relative to Smart Persistence across the scored targets. Ramp NMAE measures error during rapid changes in observed power. A gain in overall skill therefore does not compensate for worse ramp accuracy. The formulas are given in Section~\ref{sec:metric-definitions}.

Models that produce quantiles also report the CRPS approximation, 80\% interval coverage, and quantile expected calibration error. These metrics use the same valid-daylight mask and plant weighting.

Training variation is reported alongside these scores. Learned methods use fixed random seeds for data ordering and parameter initialization.

For multimodal models, a score advantage alone does not show that the trained forecast uses images. Each trained checkpoint is therefore evaluated with vision active and again with its visual embeddings set to zero. The numerical embeddings and token positions remain in place. Both passes use the same weights, forecast origins, target masks, ramp thresholds, and aggregation. Their differences in normalized error and ramp NMAE measure sensitivity to visual input under that checkpoint. This comparison keeps the training route, capacity, and regularization fixed. A separately trained numerical model would change those factors as well. A negligible vision-on/off difference can therefore coexist with a nominal advantage over a numerical model. Higher error with vision disabled indicates a positive contribution from the visual pathway under this checkpoint and intervention.

The vision-on/off pass cannot identify which part of the clip matters. The cached-latent-order, stale-image, and wrong-plant-image substitution interventions are specified in Chapter~\ref{cha:ablations}. They use the trained weights and, where possible, the same forecast origins. They are reported apart from standard model comparisons because they change the input to an existing model. A wrong-plant-image substitution needs a clip from a different test plant at the same time. Shuffled test batches allow the evaluator to find such a donor when series-major batches would contain only one plant.

Each reported result identifies the dataset track and plant partition, physical window lengths, input modalities, exact covariate fields, and forecast-origin rules. It also identifies the random seeds, numbers of windows and scored target steps, checkpoint-selection metric, tuning budget, and whether the comparison was prespecified or exploratory.
